% Pulsar system model — included by pulsar.tex §"System model"

\section*{System model}

% TODO. Cover:
%
% Network model
%   - synchronous vs asynchronous (Pulsar assumes synchronous broadcast
%     within the round window; partially-synchronous is future work)
%   - broadcast channel: authenticated, eventual delivery
%   - point-to-point channel: authenticated, encrypted
%   - public bulletin: authenticated, append-only (consensus-anchored)
%
% Setup assumptions
%   - PKI: each party has a long-term verification key advertised on the
%     public bulletin (e.g. lux/consensus identity record per HIP-0077)
%   - Random beacon: shared randomness oracle for reshare quorum selection
%     (lux consensus randomness layer)
%
% Trust model
%   - Up to t-1 corrupted parties (static), t-1 per epoch (mobile)
%   - Honest majority of (n - t + 1) needed for liveness
%
% Abort model
%   - Identifiable abort: protocol surfaces a signed complaint when a
%     deviation is detected, with the deviating party ID
%   - Slashing: governance layer consumes complaints (out of scope here)
%
% Preprocessing
%   - DKG runs once per validator-set epoch
%   - Reshare runs at epoch boundary (typical cadence: hours to days)
%   - Sign runs per-message; preprocessing-free if combined with non-cached
%     sampler, or with optional Round-0 mask precomputation per Pulsar's
%     offline/online split
%
% Rotation
%   - Reshare protocol updates shares without changing pk (per HJKY97)
%   - Members may join / leave between epochs subject to (n', t') constraint
%
% Monitoring
%   - Public artifacts: pk, commitments, signed complaints, signature outputs
%   - Private artifacts: shares, blinds, masks, samples — never logged
%
% Deployment surface
%   - Reference assumes pure software
%   - HSM integration is future work (out of scope for round 1 MPTC)
